(a) The partial derivatives of $x^2-x^4-y^4$ vanish on the curve only at $(0,0)$. Its tangent cone is $x^2=0$. Over $\mathbb{R}$, its two local branches satisfy $x=\pm y^2+O(y^6)$, so this is the tacnode in Figure 4.

(b) The singularity equations are $y=6x^5$, $x=6y^5$. They imply $xy=6x^6=6y^6$, while $xy=x^6+y^6$, hence $4xy=0$. Thus the origin is the only singular point. The tangent cone is $xy=0$, and the two real branches satisfy $y=x^5+O(x^{29})$ and $x=y^5+O(y^{29})$. This is the node.

(c) The singularity equations are $x^2(3-4x)=0$ and $2y(1+2y^2)=0$. At the origin the lowest terms are $x^3-y^2$, giving the cusp in Figure 4; over $\mathbb{R}$ its branches satisfy $y=\pm x^{3/2}(1+O(x))$ for $x\geq0$. A singular point other than the origin must satisfy $x=3/4$ and $y^2=-1/2$. Substitution in the equation gives $91=0$. Thus there are two additional singular points in characteristics $7$ and $13$, and none in other characteristics different from $2$.

(d) At a singular point, Euler's identity applied to the degree $3$ and degree $4$ terms gives $x^4+y^4=xy(x+y)=0$. Since $\operatorname{char}k\ne2$, this forces $x=y=0$. Its tangent cone is $xy(x+y)=0$, with three distinct tangent directions, so this is the triple point.
